\documentclass[10pt,a4paper]{amsart}
\usepackage[margin=19mm]{geometry}
\usepackage{amsmath,amssymb,mathtools}
\usepackage[scr=rsfs,cal=euler]{mathalfa}
\usepackage{xfrac}
\usepackage[unicode,hidelinks,bookmarks=false]{hyperref}
\newcommand{\ie}{i.e.\ }
\newcommand{\cf}{cf.\ }
\newcommand{\resp}{respectively\ }
\newcommand{\Subsection}{\subsection}
\providecommand{\nobreakdash}{\nobreak\hbox{-}}
\setlength{\emergencystretch}{2em}
\multlinegap0pt
\usepackage{bm}
\usepackage{bbm}
\usepackage{dsfont}
\usepackage{xspace}
\usepackage{cancel}
\usepackage{mathtools}
\usepackage{amsthm}
\makeatletter
\@ifundefined{lem}{\newtheorem{lem}{Lemma}}{}
\makeatother
\def\pa{\partial}
\def\th{\theta}
\let\al=\alpha
\let\vf=\varphi
\newcommand{\beno}{\begin{eqnarray*}}
\newcommand{\ben}{\begin{eqnarray}}
\newcommand{\beq}{\begin{equation}}
\newcommand{\eeno}{\end{eqnarray*}}
\newcommand{\een}{\end{eqnarray}}
\newcommand{\eeq}{\end{equation}}
\title[Optimal Convergence Estimate of the Limit from Inverse Power]{Optimal Convergence Estimate of the Limit from Inverse Power Potential to Hard Sphere Boltzmann Equation}
\author{Zheng-Nan Hu, Jin Woo Jang, Zheng-An Yao, Yu-Long Zhou}
\date{Source: arXiv:2512.21938v1 (CC BY 4.0)}
\begin{document}
\maketitle

\noindent\emph{Source and licence.} Optimal Convergence Estimate of the Limit from Inverse Power Potential to Hard Sphere Boltzmann Equation. Version arXiv:2512.21938v1 (CC BY 4.0), distributed under the Creative Commons Attribution 4.0 International licence, as recorded in the arXiv licence metadata for this exact version and shown on the abstract page https://arxiv.org/abs/2512.21938v1. Full rights notice: licenses/arxiv-2512.21938-CC-BY-4.0.txt.\par\smallskip

\noindent\textit{Editorial scope.} Counted unit: the Lem whose source label is \texttt{lemma3.3} in the article (source0.tex lines 1350--1361, source label \texttt{lemma3.3}), delivered together with the article's own complete original proof of it (source0.tex lines 1362--1525, one \texttt{proof} environment of 164 lines). No mathematics is written, repaired or re-derived: statement and proof are the article's own text line for line. Required context retained uncounted: inverse-pl-s (lines 338--343); hard-sphere (lines 345--350); lem3.1 (lines 1274--1279); 3.1 (lines 1296--1304). Every definition, display and label of those lines is retained unchanged. Omitted from this package and not redistributed: 1--337, 344--344, 351--1273, 1280--1295, 1305--1349, 1526--1913. Nothing inside a retained excerpt is omitted or altered: every retained body is delivered exactly as the article's lines are written, so no omission receipt of any kind accompanies this package. Citations: the retained excerpts carry no citation command at all, so no bibliography entry of the article is transcribed and no reference is redistributed. Reference adaptation: none was needed and none was made; every reference inside the delivered unit resolves inside it, so no unresolved reference or citation is produced. Printed slips kept literal: no printed word, symbol, hypothesis or number of the counted statement or of its proof was altered. Layout and class: the article's own class is \texttt{amsart} with packages including amsmath, mathrsfs, amssymb, graphicx, pdfsync, mathtools, color, geometry, hyperref, colonequals, todonotes; the wrapper is the lane's pinned \texttt{amsart} editorial wrapper at 10pt on A4, which restates the theorem-like environments the retained text needs and adds the article's own macro definitions that the retained text needs (\texttt{\textbackslash al}, \texttt{\textbackslash ben}, \texttt{\textbackslash beno}, \texttt{\textbackslash beq}, \texttt{\textbackslash een}, \texttt{\textbackslash eeno}, \texttt{\textbackslash eeq}, \texttt{\textbackslash pa}, \texttt{\textbackslash th}, \texttt{\textbackslash vf}), with extra wrapper packages (bm, bbm, dsfont, xspace, cancel, mathtools). The delivered numbering is the unit's own. Assets: no retained span contains a figure environment, a graphics-inclusion command, a PDF-page inclusion, a table or a TikZ picture; no image file is copied into this package, the package declares no asset dependency and no image file is redistributed with it. Licence: version arXiv:2512.21938v1 of this article is distributed under the Creative Commons Attribution 4.0 International licence, as recorded in the arXiv licence metadata for this exact version and shown on the abstract page https://arxiv.org/abs/2512.21938v1. A full rights notice is delivered at licenses/arxiv-2512.21938-CC-BY-4.0.txt. Characters: every retained excerpt is pure ASCII, so the delivered engine, XeLaTeX with its default Latin Modern text font, reports no missing character.
\begin{equation}\label{inverse-pl-s}
\begin{cases}
\pa_t f^s=Q^s(f^s,f^s)\\
f^s(0,v)=f_{\textup{in}}(v),
\end{cases}  
\end{equation}

\begin{equation}\label{hard-sphere}
\begin{cases}
\pa_t f^0=Q^0(f^0,f^0)\\
f^0(0,v)=f_{\textup{in}}(v).
\end{cases}
\end{equation}

\begin{lem}\label{lem3.1}
    Suppose f is a solution of Boltzmann equation \eqref{inverse-pl-s} or \eqref{hard-sphere} with hard potential kernel with initial data $f_{\textup{in}}\in L^1_2\cap L\log L$. Then for $k\ge 2,\ t\ge 0$, there exists constant $C(k,|f_{\textup{in}}|_{L^1_2},|f_{\textup{in}}|_{L\log L})$, such that
    \beno
|f|_{L^1_k}\le C(k,|f_{\textup{in}}|_{L^1_2},|f_{\textup{in}}|_{L\log L})|f_{\textup{in}}|_{L^1_k}.
    \eeno
\end{lem}

\begin{lem}\label{lemma3.2}
    For any constant $k\ge2$, we have
    \beq\label{3.1}
    \begin{aligned}
&\langle v'\rangle^k-\cos^k\frac{\th}{2}\langle v\rangle^k-\sin^k\frac{\th}{2}\langle v_*\rangle^k\le C_k(\cos^{k-1}\frac{\th}{2}\langle v\rangle^{k-1}\sin\frac{\th}{2}\langle v_*\rangle+\cos\frac{\th}{2}\langle v\rangle\sin^{k-1}\frac{\th}{2}\langle v_*\rangle^{k-1}),\\
&\langle v'_*\rangle^k-\sin^k\frac{\th}{2}\langle v\rangle^k-\cos^k\frac{\th}{2}\langle v_*\rangle^k\le C_k(\cos^{k-1}\frac{\th}{2}\langle v_*\rangle^{k-1}\sin\frac{\th}{2}\langle v\rangle+\cos\frac{\th}{2}\langle v_*\rangle\sin^{k-1}\frac{\th}{2}\langle v\rangle^{k-1}).
\end{aligned}
    \eeq
\end{lem}

\begin{lem}\label{lemma3.3}
Let $f^0$ be the solution of \eqref{hard-sphere} with initial data $f_{\textup{in}}\in W^{1,1}_2\cap L^1_{3}\cap L\log L$. 
    There exists  a constant $C(|f_{\textup{in}}|_{L^1_2},|f_{\textup{in}}|_{L\log L})>0$ such that
    \ben\label{3.3}
|f^0|_{W^{1,1}_2}\le |f_{\textup{in}}|_{W^{1,1}_2}e^{C(|f_{\textup{in}}|_{L^1_2},|f_{\textup{in}}|_{L\log L})|f_{\textup{in}}|_{L^1_3} t},\quad t\ge0.
    \een
    Moreover, let $2< k\in\mathbb{R}$, suppose $f_{\textup{in}}\in W^{1,1}_{k}\cap L^1_{k+1}\cap L\log L$,
    for any $T>0$, there exists a constant $C$, depending on $k,T,|f_{\textup{in}}|_{L^1_{k+1}},|f_{\textup{in}}|_{W^{1,1}_{k}},|f_{\textup{in}}|_{L\log L}$ such that
    \beno
|f^0|_{W^{1,1}_{k}}\le C(k,T,|f_{\textup{in}}|_{L^1_{k+1}},|f_{\textup{in}}|_{W^{1,1}_{k}},|f_{\textup{in}}|_{L\log L}),\quad 0\le t\le T.
    \eeno
\end{lem}

\begin{proof}[Proof of Lemma \ref{lemma3.3}]
    By \eqref{hard-sphere}, one has for any index $|\al|=1$,
    \beno
\pa_{t}\pa^{\al}f^0=Q^0(f^0,\pa^{\al}f^0)+Q^0(\pa^{\al}f^0,f^0).
    \eeno
    Recalling the  collision operator $Q^0$, 
    $$Q^0(g,f)(v)=\int_{\mathbb{R}^3\times \mathbb{S}^2}B^0(g'_*f'-g_*f)\mathrm{d}v_*\mathrm{d}\sigma.$$
    Therefore, for any smooth test function $\vf(v)$, one has
    \beq\label{3.6}
\begin{aligned}
    \langle Q^0(g,f)+Q^0(f,g),\vf\rangle=&\int B^0(g'_*f'-g_*f)\vf\mathrm{d}V+\int B^0(g'f'_*-gf_*)\vf\mathrm{d}V\\
    = & \int B^0 g_*f(\vf_*'+\vf'-\vf_*-\vf)\mathrm{d}V.
\end{aligned}
    \eeq
    Multiplying the equation \eqref{hard-sphere} with $\textup{\textup{sgn}}(\pa^{\al}f^0)\langle v\rangle^2$, integrating over $\mathbb{R}^3$, one has
    \beq\label{3.7}
    \begin{aligned}
        \frac{\mathrm{d}}{\mathrm{d}t}|\pa^{\al}f^0|_{L^1_2}=\langle Q^0(\pa^{\al}f^{0},f^0\rangle+Q^0(f^0,\pa^{\al}f^{0}),\textup{\textup{sgn}}(\pa^{\al}f^0)\langle v\rangle^2\rangle.
    \end{aligned}
    \eeq
    Note that
    \beq\label{3.8}
    \begin{aligned}
    &\langle Q^0(\pa^{\al}f^{0},f^0\rangle+Q^0(f^0,\pa^{\al}f^{0}),\textup{\textup{sgn}}(\pa^{\al}f^0)\langle v\rangle^2\rangle\\
    &\text{taking $g= f^0,f=\pa^\al f^0,\vf=\textup{\textup{sgn}}(\pa^\al f^0)\langle v\rangle^2$ in \eqref{3.6}}\\
    =&\int B^0f^0_*\pa^\al f^0(\textup{\textup{sgn}}(\pa^\al f^0)'_*\langle v'_*\rangle^2+\textup{sgn}(\pa^\al f^0)'\langle v'\rangle^2-\textup{sgn}(\pa^\al f^0)_*\langle v_*\rangle^2-\textup{sgn}(\pa^\al f^0)\langle v\rangle^2)\mathrm{d}V\\
    \le& \int B^0|\pa^\al f^0|f^0_*(\langle v'_*\rangle^2+\langle v'\rangle^2+\langle v_*\rangle^2-\langle v\rangle^2)\mathrm{d}V\\
    =& \int 2B^0|\pa^{\al}f^{0}|f^0_*\langle v_*\rangle^2\mathrm{d}V=2\pi\int_{\mathbb{R}^3\times\mathbb{R}^3}|v-v_*|f^0_*|\pa^{\al}f^0|\langle v_*\rangle^2\mathrm{d}v\mathrm{d}v_*.
    \end{aligned}
    \eeq
    By \eqref{3.7} and \eqref{3.8},
    \begin{equation}\label{3.9}
    \begin{aligned}
    \frac{\mathrm{d}}{\mathrm{d}t}|\pa^{\al}f^0|_{L^1_2}&\le 2\pi \int_{\mathbb{R}^3\times\mathbb{R}^3}|v-v_*||\pa^{\al}f^{0}|f^0_*\langle v_*\rangle^2\mathrm{d}v\mathrm{d}v_*\\
    &\le 2\pi\Big(|\pa^{\al}f^{0}|_{L^1_1}|f^0|_{L^1_2}+|\pa^{\al}f^{0}|_{L^1}|f^0|_{L^1_3} \Big)\\
    &\le 4\pi |\pa^{\al}f^{0}|_{L^1_2}|f^0|_{L^1_3}.
    \end{aligned}
    \end{equation}
    By Lemma \ref{lem3.1}, there exists a constant $C$ depending on $|f_{\textup{in}}|_{L^1_2},|f_{\textup{in}}|_{L\log L}$, such that for $t\ge0$,
    \ben\label{3.10}
|f^0|_{L^1_3}\le C(|f_{\textup{in}}|_{L^1_2},|f_{\textup{in}}|_{L\log L})|f_{\textup{in}}|_{L^1_3}.
    \een
    Combining \eqref{3.7}, \eqref{3.9} and \eqref{3.10}, we get
    \beno
\frac{\mathrm{d}}{\mathrm{d}t}|\pa^{\al}f^{0}|_{L^1_2}\le C(|f_{\textup{in}}|_{L^1_2},|f_{\textup{in}}|_{L\log L})|\pa^{\al}f^{0}|_{L^1_2}|f_{\textup{in}}|_{L^1_3}.
    \eeno
    Thanks to Gronwall's inequality, we derive for $t\ge 0$,
    \ben\label{3.12}
|\pa^{\al}f^{0}|_{L^1_2}\le |\pa^{\al}f_{\textup{in}}|_{L^1_2}e^{C(|f_{\textup{in}}|_{L^1_2},|f_{\textup{in}}|_{L\log L})|f_{\textup{in}}|_{L^1_3}t}.
    \een
    From \eqref{3.12}, 
    \beno
|f^0|_{W^{1,1}_2}= |f^0|_{L^1_2}+ \sum\limits_{|\al|=1}|\pa^{\al}f^{0}|_{L^1_2}\le |f_{\textup{in}}|_{W^{1,1}_2}e^{C(|f_{\textup{in}}|_{L^1_2},|f_{\textup{in}}|_{L\log L})|f_{\textup{in}}|_{L^1_3}t}.
    \eeno
    This completes the proof of \eqref{3.3}.\\ 
    Taking place $\textup{sgn}(\pa^{\al}f^0)\langle v\rangle^2$ with $\textup{sgn}(\pa^{\al}f^0)\langle v\rangle^{k}(2< k)$ and repeating similar argument, one has
    \ben\label{3.14}
    \frac{\mathrm{d}}{\mathrm{d}t}|\pa ^{\al}f^0|_{L^1_{k}}\le\int B^0|\pa^{\al}f^{0}|f^0_*(\langle v'_*\rangle^{k}+\langle v'\rangle^{k}+\langle v_*\rangle^{k}-\langle v\rangle^{k})\mathrm{d}V.
    \een
    
   
    By \eqref{3.1}, it gives
    \beq\label{3.15}
    \begin{aligned}
        &\langle v_*'\rangle^{k}+\langle v'\rangle^{k}+\langle v_*\rangle^{k}-\langle v\rangle^{k}\\
        \le& (\cos^{k}\frac{\th}{2}+\sin^{k}\frac{\th}{2}-1)\langle v\rangle^{k}+C_k(\cos^{k}\frac{\th}{2}+\sin^{k}\frac{\th}{2}+1)\langle v_*\rangle^{k}
        \\
        +&C_k(\langle v_*\rangle^{k-1}\langle v\rangle+\langle v_*\rangle\langle v\rangle^{k-1}).
    \end{aligned}
    \eeq
    From \eqref{3.14} and \eqref{3.15}, it reduces to estimate
    $$\int\limits_{\mathbb{R}^3\times\mathbb{R}^3\times\mathbb{S}^2}B^0|\pa^{\al}f^{0}|f^0_*(A_1+A_2+A_3)\mathrm{d}V,$$
    where 
    $$A_1=(\cos^{k}\frac{\th}{2}+\sin^{k}\frac{\th}{2}-1)\langle v\rangle^{k},\quad A_2=C_k(\cos^{k}\frac{\th}{2}+\sin^{k}\frac{\th}{2}+1)\langle v_*\rangle^{k},$$
    $$A_3=C_k(\langle v_*\rangle^{k-1}\langle v\rangle+\langle v_*\rangle\langle v\rangle^{k-1}).$$
   
  \textbf{Estimate of $A_1$.} First,
    \ben\label{3.16}
    \int B^0|\pa^{\al}f^0|f^0_* A_1 \mathrm{d}V=-C_1(k)\int_{\mathbb{R}^3\times\mathbb{R}^3}|v-v_*||\pa^{\al}f^0|f^0_*\langle v\rangle^{k}\mathrm{d}v\mathrm{d}v_*,
    \een
    where $C_1(k)=\int_{\mathbb{S}^2}\frac{1}{4}(1-\sin^{k}\frac{\th}{2}-\cos^{k}\frac{\th}{2})\mathrm{d}\sigma>0$ depending only on $k$.
 Let
   $\nu(v)=
\int_{\mathbb{R}^3}|v-v_*|f^0(v_*)\mathrm{d}v_*$. On the one hand, for any $K>0$,
\beq\label{3.17}
\begin{aligned}
    \nu(v)&\ge \int_{\mathbb{R}^3,|v-v_*|\ge K}|v-v_*|f^0(v_*)\mathrm{d}v_*\\
    &\ge K\int_{\mathbb{R}^3,|v-v_*|\ge K}f^0(v_*)\mathrm{d}v_*=K(|f^0|_{L^1}-\int_{\mathbb{R}^3,|v-v_*|\le K}f^0(v_*)\mathrm{d}v_*).
\end{aligned}
\eeq
By mass conservation, one has $|f^0|_{L^1}=|f_{\textup{in}}|_{L^1}$. Note that for any $R>1$,
\beq\label{estimate of convolution-1}
\begin{aligned}
\int_{\mathbb{R}^3,|v-v_*|\le K}&f^0(v_*)\mathrm{d}v_*\le \int_{|f^0|\ge R,|v-v_*|\le K}f^0(v_*)\mathrm{d}v_*+\int_{|f^0|\le R,|v-v_*|\le K}f^0(v_*)\mathrm{d}v_*\\
\le & \int_{|f^0|\ge R,|v-v_*|\le K}f^0(v_*)\log f^0(v_*)\mathrm{d}v_* (\log R)^{-1}+ \frac{4}{3}\pi K^3R\\
\le & \int_{|f^0|\ge 1}f^0(v_*)\log f^0(v_*)\mathrm{d}v_* (\log R)^{-1}+ \frac{4}{3}\pi K^3R\\
\le & (H(f^0)-\int_{|f^0|\le 1}f^0(v_*)\log f^0(v_*)\mathrm{d}v_*)(\log R)^{-1}+\frac{4}{3}\pi K^3R.
\end{aligned}
\eeq
By the well-known Boltzmann H-theorem, one has $H(f^0)\le H(f_{\textup{in}})$. Using a basic inequality,
$$x\log y-x\log x\le y,\quad y>0,x\ge 0.$$
One can take $x=f^0(v_*),y=e^{-|v_*|^2}$ to get
$$-\int_{|f^0|\le 1}f^0(v_*)\log f^0(v_*)\mathrm{d}v_*\le \int_{\mathbb{R}^3} f^0(v_*)|v_*|^2\mathrm{d}v_*+\int_{\mathbb{R}^3} e^{-|v_*|^2}\mathrm{d}v_*\le |f_{\textup{in}}|_{L^1_2}+\pi^{3/2}.$$
Therefore, \eqref{estimate of convolution-1} gives
\ben\label{estimate of convolution-2}
\int_{\mathbb{R}^3,|v-v_*|\le K}f^0(v_*)\mathrm{d}v_*\le (\log R)^{-1}(H(f_{\textup{in}})+|f_{\textup{in}}|_{L^1_2}+\pi^{3/2})+\frac{4}{3}\pi K^3R.
\een
From \eqref{3.17} and \eqref{estimate of convolution-2},
\beno
\nu(v)\ge K\Big(|f_{\textup{in}}|_{L^1}-(\log R)^{-1}(H(f_{\textup{in}})+|f_{\textup{in}}|_{L^1_2}+\pi^{3/2})-\frac{4}{3}\pi K^3R\Big).
\eeno
 By the assumption $f_{\textup{in}}\in L^1_{2}\cap L\log L$, one has $H(f_{\textup{in}})+|f_{\textup{in}}|_{L^1_{2}}<\infty$. We can take $R=\exp(\frac{4(H(f_{\textup{in}})+\pi^{3/2}+|f_{\textup{in}}|_{L^1_{2}})}{|f_{\textup{in}}|_{L^1}})$ and $K=(\frac{3}{16\pi}R^{-1}|f_{\textup{in}}|_{L^1})^{1/3}$. Then
\ben\label{3.21}
\nu(v)\ge \frac{K}{2}|f_{\textup{in}}|_{L^1}>0.
\een
On the other hand, it is easy to see
\beq\label{3.22}
\nu(v)\ge |v||f^0|_{L^1}-|f^0|_{L^1_1}\ge |v||f_{\textup{in}}|_{L^1}-|f_{\textup{in}}|_{L^1_2}.
\eeq
    By \eqref{3.21} and \eqref{3.22}, there exists a constant $C_2$ depending on $|f_{\textup{in}}|_{L^1},|f_{\textup{in}}|_{L^1_2},|f_{\textup{in}}|_{L\log L}$, such that 
    \ben\label{3.23}
\nu(v)\ge C_2\langle v \rangle .
    \een
    From \eqref{3.16} and \eqref{3.23},
\ben\label{3.24}
    \int B^0|\pa^{\al}f^0|f^0_*A_1\mathrm{d}V\le -C_1(k)C_2|\pa^{\al}f^0|_{L^1_{k+1}}.
\een

 \textbf{Estimate of $A_2$.} By Young's inequality,
$$
\begin{aligned}
\int B^0|\pa^{\al}f^0|f^0_* A_2\mathrm{d}V&= C_k\int_{\mathbb{S}^2}\frac{1}{4}(\cos^{k}\frac{\th}{2}+\sin^{k}\frac{\th}{2}+1)\mathrm{d}\sigma\int |v-v_*||\pa^{\al}f^0|f^0_*\langle v_*\rangle^{k}\mathrm{d}v\mathrm{d}v_*\\
&\le C_k\int_{\mathbb{R}^3\times\mathbb{R}^3}|\pa^{\al}f^0|f^0_*\langle v\rangle\langle v_*\rangle^{k}\mathrm{d}v\mathrm{d}v_*+
C_k\int_{\mathbb{R}^3\times\mathbb{R}^3}|\pa^{\al}f^0|f^0_*\langle v_*\rangle^{k+1}\mathrm{d}v\mathrm{d}v_*\\
&\le C_k\epsilon |\pa^{\al}f^0|_{L^1_{k+1}}|f^0|_{L^1}+C_k(\epsilon^{-\frac{1}{k}}+1)|\pa^{\al}f^0|_{L^1}|f^0|_{L^1_{k+1}}.
\end{aligned}
$$
Therefore,
\ben\label{3.25}
\int B^0|\pa^{\al}f^0|f^0_* A_2\mathrm{d}V\le C_k\epsilon |\pa^{\al}f^0|_{L^1_{k+1}}|f_{\textup{in}}|_{L^1}+C(k,\epsilon)|\pa^{\al}f^0|_{L^1}|f_{\textup{in}}|_{L^1_{k+1}}.
\een

 \textbf{Estimate of $A_3$.} By Young's inequality,
\beq\label{3.26}
\begin{aligned}
   & \int B^0|\pa^{\al}f^0|f^0_* A_3\mathrm{d}V=C_k\int B^0|\pa^{\al}f^0|f^0_*(\langle v\rangle\langle v_*\rangle^{k-1}+\langle v_*\rangle\langle v\rangle^{k-1})\mathrm{d}V\\
    &\le C_k\Big(\epsilon\int_{\mathbb{R}^3\times\mathbb{R}^3} |v-v_*||\pa^{\al}f^0|f^0_*\langle v\rangle^{k}\mathrm{d}v\mathrm{d}v_*+(\epsilon^{-\frac{1}{k-1}}+\epsilon^{-\frac{k}{k-1}})\int_{\mathbb{R}^3\times\mathbb{R}^3} |v-v_*||\pa^{\al}f^0|f^0_*\langle v_*\rangle^{k}\mathrm{d}v\mathrm{d}v_*\Big)\\
    &\le C_k\epsilon |\pa^{\al}f^0|_{L^1_{k+1}}|f_{\textup{in}}|_{L^1_2}+C_k(\epsilon^{-\frac{1}{k-1}}+\epsilon^{-\frac{k}{k-1}})|\pa^{\al}f^0|_{L^1_2}|f_{\textup{in}}|_{L^1_{k+1}}.
\end{aligned}
\eeq
Combining \eqref{3.24}, \eqref{3.25} and \eqref{3.26}, taking $\epsilon$ small enough, we derive
$$
\int B^0|\pa^{\al}f^0|f^0_*(A_1+A_2+A_3)\mathrm{d}V\le -\frac{1}{2}C_1(k)C_2|\pa^{\al}f^0|_{L^1_{k+1}}+C_3(k,|f_{\textup{in}}|_{L^1_2},|f_{\textup{in}|_{L\log L}})|\pa^{\al}f^0|_{L^1_2}|f_{\textup{in}}|_{L^1_{k+1}}.
$$
By \eqref{3.12},
\ben\label{3.27}
\frac{\mathrm{d}}{\mathrm{d}t}|\pa^{\al}f^0|_{L^1_{k}}\le -\frac{1}{2}C_1(k)C_2|\pa^{\al}f^0|_{L^1_{k+1}}+C_3|\pa^{\al}f_{\textup{in}}|_{L^1_2}e^{C|f_{\textup{in}}|_{L^1_3}t}|f_{\textup{in}}|_{L^1_{k+1}}.
\een
Solving \eqref{3.27},
\ben\label{3.28}
|\pa^{\al}f^0|_{L^1_{k}}\le |\pa^{\al}f_{\textup{in}}|_{L^1_{k}}\exp\{-\frac{1}{2}C_1C_2t\}+\frac{C_3|\pa^{\al}f_{\textup{in}}|_{L^1_{2}}|f_{\textup{in}}|_{L^1_{k+1}}}{C|f_{\textup{in}}|_{L^1_3}+\frac{1}{2}C_1C_2}\Big(\exp\{C|f_{\textup{in}}|_{L^1_3}t\}-\exp\{-\frac{1}{2}C_1C_2t\}\Big).
\een
For any $T>0$, \eqref{3.28} gives a bound of $|f^0|_{W^{1,1}_{k}}$, which depends on $k$, $T$, $|f_{\textup{in}}|_{W^{1,1}_{k}}$ , $|f_{\textup{in}}|_{L^1_{k+1}}$ and $|f_{\textup{in}}|_{L\log L}$.
\end{proof}

\end{document}
